% ---------------------------------------------------------------------------
% H3X - Embedded Controls Intern
% Tailored from bases/embedded_controls.tex on 2026-08-26
% Source context: recruiter outreach saved in job_description.txt
%
% The company name belongs in the filename only, not the rendered document.
% ---------------------------------------------------------------------------
\documentclass[10pt,letterpaper]{article}
\usepackage{resume}

% Preserve the reference resume's one-page density without reducing type size.
\setlist[itemize]{topsep=1pt,itemsep=0.5pt,after=\vspace{1pt}}
\renewcommand{\entryspace}{\vspace{0pt}}

% Use vertical bars instead of hyphens as contact separators.
\renewcommand{\resumeheader}{%
  \name{Akash Karthik}
  \contact{San Jose, CA \textbar{} akashkarthik473@gmail.com \textbar{} linkedin.com/in/akash-karthik473 \textbar{} github.com/akashkarthik473 \textbar{} 9254069507}
}

\begin{document}
\bodyfont

\resumeheader
\educationsection

\sectionheader{TECHNICAL SKILLS}

\techline{\textbf{Languages}: C/C++, Python, Assembly, Java}
\techline{\textbf{Electric Powertrain Controls}: Fixed-point PI/PID control, power limiting, energy management, saturation, anti-windup, slew-rate limiting}
\techline{\textbf{Embedded Systems}: STM32H755, TTC 60 VCU, PX4, real-time firmware, CAN, SPI, microcontroller I/O}
\techline{\textbf{Test \& Tools}: Software-in-the-loop, CAN trace analysis, TASKING, BusMaster, MoTeC i2 Pro, GDB, Git}

\sectionheader{EXPERIENCE}

\roleline{Software Lead}{Jun 2026 to Present}
\orgline{\spartanracing}
\begin{itemize}
  \item Lead the software organization for a Formula SAE electric racecar, directing \textbf{12 engineers} across \textbf{10+ concurrent projects} in four-motor controls, BMS and dashboard firmware, thermal derating, plant modeling, and telemetry
  \item Direct four-motor controls and validation, defining a phased torque vectoring and traction control program and coordinating battery, motor, and inverter plant models for software- and hardware-in-the-loop testing
\end{itemize}

\entryspace

\roleline{Software Engineer Intern}{Feb 2026 to Present}
\orgline{Argus Defense \textbar{} San Jose, CA}
\begin{itemize}
  \item Tuned PID flight-control loops on a PX4 flight controller to reduce vertical and horizontal position drift in autonomous flight
  \item Integrated and tuned visual-inertial odometry for onboard state estimation, improving position-hold accuracy
  \item Developed and validated embedded flight-control firmware in C on PX4, iterating control gains against logged flight telemetry
\end{itemize}

\entryspace

\roleline{Software Controls Engineer}{Jul 2025 to Jun 2026}
\orgline{\spartanracing}
\begin{itemize}
  \item Designed a lap-adaptive energy-management algorithm that sets a power setpoint from cumulative energy versus budget, with measured power tracking the commanded setpoint within \textbf{0.01\%}
  \item Developed embedded power-limiting firmware with a PI controller on the VCU, maximizing motor output while enforcing the Formula SAE \textbf{80 kW} limit
  \item Built CAN diagnostics firmware for live power and torque frames, status bits, and error counters for real-time BusMaster monitoring
\end{itemize}

\entryspace

\roleline{Software Intern}{Aug 2024 to Jun 2025}
\orgline{\spartanracing}
\begin{itemize}
  \item Built deterministic entry and exit state logic for power limiting from measured power, requested torque, and threshold flags
  \item Tuned the PI torque controller used for power limiting to within \textbf{2\%} error, smoothing torque transitions to preserve drivability
  \item \textbf{Results:} 1st in endurance at MIS EV 2025
\end{itemize}

\sectionheader{PROJECT EXPERIENCE}

\roleline{Formula SAE EV Vehicle Simulation}{Jul 2025}
\orgline{\spartanracing}
\techline{Python, Pygame, Matplotlib, PID control, vehicle dynamics, CSV logging}
\begin{itemize}
  \item Built a physics-based EV simulation to model acceleration from vehicle mass, drag, gearing, and motor torque
  \item Implemented PID-based power limiting to the \textbf{80 kW} rule and logged runs to CSV for repeatable tuning and validation
\end{itemize}

\entryspace

\roleline{SR-Wireless-CAN, Vehicle Telemetry Platform}{2026}
\orgline{\spartanracing \textbar{} Backend lead}
\techline{Python, python-can, FastAPI, WebSockets, SQLite, Docker}
\begin{itemize}
  \item Proposed and implemented wireless firmware flashing for the VCU by decoding the vendor tool's CAN sequence from trace captures and reimplementing it in Python
  \item Developed CAN capture and diagnostics tooling with DBC-driven signal decoding for live vehicle testing and recorded trace replay
\end{itemize}

\end{document}
